\section{Background and related work}
\label{sec:related}

\paragraph{Metastable failures.} Bronson et al.\ \cite{bronson2021metastable} define
the class through its three elements: a trigger, a sustaining
effect, and a degraded stable state. Huang et al.\ \cite{huang2022wild}
study production occurrences, including a look-aside-cache
reproduction in which a cache-hit-rate drop sustains overload.
Farahbakhsh et al.\ state the prediction agenda - a model sketch for
whether a given request-response structure admits metastable states
\cite{farahbakhsh2025modeling}, with a causal characterization in
\cite{farahbakhsh2026characterizing}. These works establish the class and
the agenda. This paper exhibits a new sustaining mechanism,
hit-rate-dependent KV write amplification through full-context
re-prefill, with a measured probability band, measured dependence on
tier capacity and router span, and a calibrated model. The mechanism
differs from the look-aside instance in where the amplification
lands: there, misses amplify load onto a backing service behind the
cache; here, a miss writes its re-prefilled context back into the
same finite pool it missed in - write amplification into the cache's
own capacity. Metronome \cite{meng2026metronome} applies the metastable
label to LLM KV memory in a setting where full-duplex sessions pin
monotonically growing per-session state and a linear open-loop fill
model predicts an abrupt stall; that model has one absorbing state
and no feedback, hit-rate, hysteresis, or recovery structure. Its
14-of-20 versus 6-of-20 outcome split under identical conditions
independently corroborates outcome bimodality at fixed operating
points. The metastability here is of the Bronson kind: a sustaining
loop, demonstrated by hysteresis.

\paragraph{Bifurcation treatments of AI serving loops.} van Rooyen
\cite{vanrooyen2026collapse} develops first-order collapse, closed-form
spinodals, cusp-organized bistability, and hysteretic recovery for
an AI system's operating loop under congestion-dependent feedback.
This is the closest published vocabulary match to the equilibrium
structure here, but the physical feedback differs: there,
uncertified output re-enters the loop as added load (retry-type
amplification); here the feedback channel is the cache state itself.
On abstract-level review (full-text review precedes submission,
Section~\ref{sec:limitations}), that work carries no cache state variable, no phase
boundary in (capacity, load), no workload calibration, and no
serving-system validation of a cache mechanism. The two are
complementary instantiations of the same fold-and-hysteresis
mathematics; ours is the cache-feedback instantiation with
measured-constant calibration and multi-replica validation.

\paragraph{Queueing stability for LLM inference.} Nie et al.\ \cite{nie2026queueing}
give a queueing-theoretic stability framework for LLM inference
under KV memory constraints; Dai et al.\ \cite{dai2025throughput} prove
throughput-optimality of work-conserving schedulers in a fluid-limit
framework covering agentic workloads; Dong and Cao \cite{dong2026flow}
derive stability conditions for admission-budgeted scheduling. In
all of these the service rate is cache-independent and memory enters
as a constraint. With hit-rate feedback the effective service rate
is state-dependent, and the stability region acquires a coexistence
band that work-conservation arguments do not see. Our closed-loop
measurements (Section~\ref{sec:closedloop}) locate the arrival semantics under which
the constraint-style analyses apply. Ao et al.\ analyze
service-induced congestion in memory-constrained serving
\cite{ao2026congestion}, in a line with fluid-guided scheduling under
endogenous per-request memory growth \cite{ao2025fluid}; its mechanism is
memory growth during service, not prefix-cache feedback (full-text
review pending, Section~\ref{sec:limitations}).

\paragraph{Cache-aware routing.} DualMap \cite{yuan2026dualmap} engineers the
affinity-versus-load trade-off with dual-hash power-of-two-choices
mapping and SLO-triggered fallback. Continuum \cite{li2025continuum} pins
KV across tool-call gaps with a TTL. CacheScout
\cite{zhang2026cachescout} learns agent execution transitions to guide
eviction and prefetch. These works optimize the trade-off's
performance. We show the trade-off has a stability dimension: the
affinity term keeps the fleet on the warm branch and carries the
recovery path (Section~\ref{sec:affinity}, one measured run), and the span of the
balancer sets collapse probability and blast radius - quantities
absent from their evaluations.

\paragraph{Tiered KV caches and restore bandwidth.} Mooncake \cite{qin2025mooncake}
trades storage for computation in a KVCache-centric architecture and
motivates early rejection under overload. CacheFlow
\cite{nian2026cacheflow} identifies KV restoration as a dominant
bottleneck and parallelizes it. Kareto \cite{zheng2026kareto} treats tier
sizing as a multi-objective steady-state optimization. Unclaimed in
this literature is the persistent congested serving regime pinned at
restore throughput, with persistence set by tier capacity (Sections
\ref{sec:congestion}, \ref{sec:sizeseries}), and the stability discontinuity in tier sizing - the
fall-through to cold collapse at insufficient capacity - which a
steady-state Pareto analysis cannot expose.

\paragraph{Classical cache theory.} The characteristic-time frame of Che et al.\
\cite{che2002hierarchical}, with the validity analysis of Fricker, Robert,
and Roberts \cite{fricker2012versatile}, describes an LRU cache whose
retention window is set by the aggregate write rate. The prefix
cache fits this frame with one addition outside its assumptions: the
write rate depends on the hit rate, because misses write full
contexts (Section~\ref{sec:retention}). Multiple cache steady states are themselves
classical: Rosensweig et al.\ \cite{rosensweig2013steady} show cache
networks can be non-ergodic, with the steady state depending on
initial placement. That multiplicity is combinatorial, arising from
placement dependencies across a topology at fixed demand, with no
load parameter, no capacity-load boundary, and no hysteresis sweep;
the multiplicity here comes from a scalar write-rate feedback within
a single tier, yielding a load-parameterized coexistence band and a
hysteresis loop. Denning's thrashing analysis \cite{denning1968thrashing},
with its multiprogramming-level control, is the classical ancestor
of the operational response to such regimes; the differences here
are a measured probability band, hysteresis, and coordination-layer
dependence rather than a qualitative regime.

\paragraph{Observation base.} CONCUR \cite{chen2026concur} is the closest published
statement of the phenomenon: in agentic batch inference, KV usage
pins at 80--100\% while the prefix-cache hit rate collapses and stays
low for most of execution (its Figure~3, from a large-scale
deployment), with a measured hit-rate knee against concurrency,
mitigated by an AIMD admission controller keyed on cache usage and
hit rate. It stops at the observation and a preventive
fixed-constant heuristic: hit rate is a measured signal there, never
a modeled state variable; no absorbing state, hysteresis,
probability band, or recovery semantics is measured or claimed; and
its mechanism attribution is exogenous (agents paused for tool calls
losing LRU recency), not the miss-write amplification that closes
the loop here. TraceLab \cite{zhu2026tracelab} characterizes coding-agent
traffic on a public corpus, including a retention-window-to-hit-rate
sweep and the finding that idle gaps beyond $\sim$5 minutes drive most
misses - the gap structure our protocol's gap cap controls
(Section~\ref{sec:workload}). Yuan et al.\ \cite{yuan2026agentic} characterize agentic
workloads and, per abstract-level review, warn of a thrashing regime
once aggregate agent context exceeds GPU memory. Ranganathan et al.\
\cite{ranganathan2025incidents} taxonomize 156 high-severity production
LLM inference incidents; Nixon et al.\ \cite{nixon2026yearserving} release
a year-long production serving trace. The phenomenon class is widely
observed and, prior to this work, unmodeled for the prefix-cache
feedback channel.
